\uuid{UInk}
\titre{Inégalité}
\auteur{bodin ; hanani ; mzari}
\organisation{exo7}
\licence{CC BY-NC-SA 4.0}
\datecreate{2020-01-20}
\niveau{L1}
\module{Analyse}
\chapitre{Propriétés de R}
\sousChapitre{Propriétés des nombres réels}
\difficulte{1}
\theme{Propriétés de nombres réels, inégalité}
\exo7qcmid{lille-1.reels.008}

\contenu{
\begin{question}
\qclassification{120.03}
\qtags{motcle=inégalité}
Soient $x,y \in \Rr$ tels que $x \le 2y$.
Quelles sont les assertions vraies ?
\begin{answers}
    \bad{$x^2 \le 2xy$}

    \bad{$y\le \frac{x}{2}$}

    \good{$2x \le x+2y$}

    \good{$-2y \le -x$}
\end{answers}
\begin{explanations}
Lorsque l'on multiplie par un nombre négatif alors le sens de l'inégalité change. Il faut faire attention lorsque l'on multiplie par $x$, car le sens de l'inégalité est changé ou pas selon que $x$ soit négatif ou positif !
\end{explanations}
\end{question}
}
